\clearpage\SourceLeaf{151}%
\pdfbookmark[1]{Gauss to von Lindenau}{lindenau1808aug}
\fancyhead[L]{\small GAUSS TO VON LINDENAU}
\begin{center}
{\small GAUSS TO VON LINDENAU.\par}
\medskip\hrule\medskip
{\small Monatliche Correspondenz zur Beförderung der Erd- und Himmels-Kunde,\par
edited by Baron \textsc{von Zach}. September 1808.\par}
\medskip\hrule
\end{center}\medskip
\begin{flushright}Göttingen, 30 August 1808.\end{flushright}
\begin{center}\small
—\ —\ —\ —\ —\ —\ —\ —\ —\ —\ —\ —\ —\ —\ —\ —\ —
\end{center}
I have the honour of sending you a paper for the \emph{Mon. Corresp.} [p.\ 141 of this volume] on a problem whose practical application deserves to be strongly recommended. It seems to me that this method of determining the elevation of the pole allows very high precision in practice and gives the most reliable results when there is no opportunity to observe with repeating circles. At least it avoids all the circumstances that can sometimes make the results of the usual method using solar altitudes doubtful. These are the reading of the instrument, errors of the sextant, dark glasses, the glass roof, and refraction. Admittedly, each combination of three stars always gives only one result. But however many observations one makes of the meridian altitudes of the Sun, the errors due to the circumstances I have mentioned, apart from the first, always remain in full force. And there is nothing to prevent one from repeating the same observations on several nights or even making several combinations in a single night, provided the combinations are chosen so that the three observations follow one another at intervals that are not too long. Last night I observed the three stars listed in the example in the paper once again at
\SourceLeaf{152}%
\expandafter\def\csname HistoricalNoteText7\endcsname{Error of the sextant at $105^0\,52'= -49''4$.}%
a somewhat different altitude and found $51^0\,31'\,54''4$\HistoricalNote{7} for the elevation of the pole. The three results are as follows:
\begin{center}
\begin{tabular}{@{}r@{\ }r@{\quad}c@{\quad}l@{}}
August & 25 & $.$ & $51^0\quad31'\quad56''7$\\
       & 27 & $.$ & $51\quad31\quad51.5$\\
       & 30 & $.$ & $51\quad31\quad54.4$\\\cline{2-4}
\multicolumn{2}{@{}r@{\quad}}{mean} & $.$ & $51\quad31\quad54.2$
\end{tabular}
\end{center}
exactly the value \textsc{Mayer} found for our elevation of the pole. Although this degree of agreement must in part be considered accidental in this instance, I nevertheless believe that, by repeating the observations more often and taking more of them with well-determined stars, one can always come within a few seconds of the true elevation of the pole. Everything depends solely on the magnification and clarity of the telescope and the observer's care. From meridian altitudes of the Sun that were always in excellent agreement with one another, I had obtained $51^0\,32'\,32''$ several times with a sextant at the time of the solstice, so the result was in error by more than half a minute. I have yet to examine closely enough whether my sextant's error comes solely from the entire arc being too large. The very method I recommend will, however, make that very easy to check in future.

\par\bigskip\centerline{\rule{.27\linewidth}{.4pt}}
